\begin{titlepage}
  \begin{center}

  {\Huge system builder library}

  \vspace{25mm}

  \includegraphics[width=0.90\textwidth,height=\textheight,keepaspectratio]{img/AFRL.png}

  \vspace{25mm}

  \today

  \vspace{15mm}

  {\Large Jay Convertino}

  \end{center}
\end{titlepage}

\tableofcontents

\newpage

\section{Usage}

\subsection{Introduction}

\par
System build is a library that will parse a yaml project file, and a yaml command
file to generate a list of commands to execute in threads. These can be concurrent
or singular in fashion.

\subsection{Dependencies}

\par
The following are the dependencies of the cores.

\begin{itemize}
  \item python 3.X
  \item python progressbar2
\end{itemize}

\subsection{Example Creator}
\par
The creator is responsable for instantiating the library, preparing the project, and launching the system builder library to build the project.
\begin{lstlisting}[language=Python]
from system.build import *

cmd_compiler = commandCompiler("projects.yml", "commands.yml")

cmd_compiler.create()

projects = cmd_compiler.getResult()

cmd_exec = commandExecutor(projects, args.dryrun)

cmd_exec.runProject()

\end{lstlisting}

\subsection{Example Commands Yaml file}
This yaml file is injested so commands are executed with certain steps for the final end product.
\begin{lstlisting}[language=XML]
fusesoc:
  cmd_0: ["__CHECK_SKIP__{_pwd}/output/hdl/{_project_name}/AFRL_project_veronica_axi_baremetal_1.0.0.bit"]
  cmd_1: ["fusesoc", "--cores-root", "{path}", "run", "--build", "--work-root", "output/hdl/{_project_name}", "--target", "{target}", "{project}"]
script:
  cmd_1: ["{exec}", "{file}", "{_project_name}", "{args}"]
gcc_riscv32:
  cmd_0: ["__CHECK_SKIP__{_pwd}/output/bin/riscv/bin/riscv32-unknown-elf-gcc"]
  cmd_1: ["make", "-C", "{_pwd}/{path}", "clean"]
  cmd_2: ["__CWD__{_pwd}/{path}", "./configure", "--prefix={_pwd}/output/bin/riscv", "--disable-linux", "--with-arch=rv32imac", "--with-abi-ilp32"]
  cmd_3: ["make", "-C", "{_pwd}/{path}", "-j", "8"]
openocd_riscv:
  cmd_0: ["__CHECK_SKIP__{_pwd}/output/bin/openocd/bin/openocd"]
  #cmd_1: ["__CWD__{_pwd}/{path}", "git", "apply", "--ignore-whitespace", "../patch/pmpregs.patch"]
  cmd_2: ["__CWD__{_pwd}/{path}", "./bootstrap"]
  cmd_3: ["__CWD__{_pwd}/{path}", "./configure", "--prefix={_pwd}/output/bin/openocd", "--enable-ftdi", "--enable-dummy", "--enable-jtag_vpi"]
  cmd_4: ["make", "-C", "{_pwd}/{path}", "-j", "8"]
  cmd_5: ["make", "install", "-C", "{_pwd}/{path}"]
fpga_baremetal_examples:
  cmd_0: ["__CHECK_SKIP__{_pwd}/output/apps"]
  cmd_1: ["mkdir", "-p", "{_pwd}/{path}/cmake"]
  cmd_2: ["__CWD__{_pwd}/{path}/cmake", "__ENV_PATH__{_pwd}/output/bin/riscv/bin/", "cmake", "../", "-DCMAKE_RUNTIME_OUTPUT_DIRECTORY={_pwd}/output/apps", "-DCMAKE_PREFIX_PATH={_pwd}/output/bin/riscv/bin/", "-DBUILD_EXAMPLES_ALL=ON", "-DCMAKE_TOOLCHAIN_FILE={_pwd}/{path}/arch/riscv/riscv.cmake"]
  cmd_3: ["__ENV_PATH__{_pwd}/output/bin/riscv/bin/", "make", "-C", "{_pwd}/{path}/cmake"]
buildroot:
  cmd_1: ["rm", "-rf", "{_pwd}/output/linux/{_project_name}"]
  cmd_2: ["make", "-C", "{path}", "distclean"]
  cmd_3: ["make", "O={_pwd}/output/linux/{_project_name}", "-C", "{path}", "{config}"]
  cmd_4: ["make", "O={_pwd}/output/linux/{_project_name}", "-C", "{path}"]
genimage:
  cmd_1: ["mkdir", "-p", "{_pwd}/output/genimage/tmp/{_project_name}"]
  cmd_2: ["genimage", "--config", "{path}/{_project_name}.cfg"]
\end{lstlisting}

\subsection{Example Projects Yaml file}
This takes the commands and fills out the information required and order of operations for creating the finaly end result project.
\begin{lstlisting}[language=XML]
zed_fmcomms2-3_linux_busybox_sdcard:
  concurrent:
    fusesoc: &fusesoc_fmcomms2-3
      path: hdl
      project: AFRL:project:fmcomms2-3:1.0.0
      target: zed_bootgen
    buildroot: &buildroot_fmcomms2-3
      path: sw/linux/buildroot-afrl
      config: zynq_zed_ad_fmcomms2-3_defconfig
  sequential:
    script: &output_files_fmcomms2-3
      exec: python
      file: py/output_gen.py
      args: "--rootfs output/linux --bootfs output/hdl --dest output/sdcard"
    genimage:
      path: img_cfg
zc702_fmcomms2-3_linux_busybox_sdcard:
  concurrent:
    fusesoc:
      <<: *fusesoc_fmcomms2-3
      target: zc702_bootgen
    buildroot:
      <<: *buildroot_fmcomms2-3
      config: zynq_zc702_ad_fmcomms2-3_defconfig
  sequential:
    script:
      <<: *output_files_fmcomms2-3
    genimage:
      path: img_cfg
\end{lstlisting}

\section{Code Documentation} \label{Code Documentation}

\par
Natural docs is used to generate documentation for this project. The next lists the following sections.

\begin{itemize}
\item \textbf{init} python init code\\
\item \textbf{builder} system builder library.\\
\item \textbf{creator} system builder example code.\\
\end{itemize}

